\section{Related Work}
\label{sec:related_work}

The literature intersecting this research encompasses two dominant trajectories: multiclass sentiment classification in code-mixed and Romanized low-resource languages, and transformer robustness under adversarial or natural orthographic perturbation.

\subsection{Multiclass Sentiment Classification in Roman Urdu}
Early computational investigations in Roman Urdu sentiment analysis relied heavily on surface-level feature engineering, employing Bag-of-Words (BoW) representations, character n-grams, and classical machine learning estimators such as Support Vector Machines (SVM), Naive Bayes, and Random Forests~\cite{chandio2022sentiment}. While character n-grams partially alleviate minor typing slips by capturing sub-word substrings, classical estimators are fundamentally incapable of modeling long-range syntactic dependencies, negation scope, and polysemy.

The introduction of multilingual pre-trained encoders catalyzed significant progress. Devlin et al.~\cite{devlin2019bert} established the efficacy of multilingual BERT (mBERT), pre-trained across 104 languages on Wikipedia corpora. Conneau et al.~\cite{conneau2020unsupervised} subsequently introduced XLM-RoBERTa (XLM-R), trained on 2.5TB of filtered CommonCrawl data in 100 languages, demonstrating substantial performance advantages in cross-lingual transfer benchmarks. 

In a notable recent study, Nazir et al.~\cite{nazir2025leveraging} proposed CMSA-mBERT for three-class sentiment analysis across 48,324 Roman Urdu and Roman Punjabi code-mixed reviews. Their model achieved a benchmark accuracy of 79.94\%, outperforming traditional ML and recurrent architectures (LSTM, BiLSTM, GRU). Crucially, the authors observed that standard preprocessing techniques (stripping stopwords, punctuation, and emojis) produced only an insubstantial 2.71\% improvement, citing unstandardized Romanized orthography as an intractable limitation. However, their experimental design was restricted to a single clean 80/20 train/test split, leaving the models' resilience to spelling noise untested.

\subsection{Vocabulary Adaptation and Embedding Stability}
Multilingual language models typically suffer from representation degradation when applied to languages absent from their pre-training corpora. Because Roman Urdu is underrepresented in mBERT and XLM-RoBERTa, native words are fragmented into highly localized subword units. 

Addressing this phenomenon, Khan et al.~\cite{khan2026romeva} introduced ROMEVA, an anchor-loss vocabulary expansion technique designed to preserve embedding-space geometry when adapting mBERT to Roman Urdu using the RUWV-NSR dataset~\cite{ahmed2024ruwvnsr}. Their findings indicated that while unconstrained naive fine-tuning disrupted the pre-trained semantic geometry more severely than constrained methods, it paradoxically achieved superior downstream classification accuracy. Nevertheless, Khan et al. focused exclusively on embedding drift rather than input-level orthographic perturbation, did not evaluate parameter-efficient tuning paradigms, and evaluated solely on unaltered review splits.

\subsection{Robustness and Perturbation Evaluation}
Assessing the brittle nature of deep neural language models under input perturbations is an active subfield. Moradi and Samwald~\cite{moradi2021evaluating} systematically evaluated transformer architectures under synthetic character-level noises (random character insertion, deletion, and swap), finding severe degradation across standard benchmarks.

However, existing perturbation benchmarks predominantly introduce synthetic typographic noise (e.g., keyboard layout typos) rather than authentic, sociolinguistically validated spelling variations. In Roman Urdu, spelling variation is not random noise but a structured phonological phenomenon. By leveraging an authentic, human-documented variation lexicon, our work provides the first controlled robustness benchmark for Roman Urdu sentiment classification.

\subsection{Parameter-Efficient Fine-Tuning (PEFT)}
As language models have scaled, parameter-efficient adaptation methods have emerged as powerful alternatives to full fine-tuning. Hu et al.~\cite{hu2022lora} introduced Low-Rank Adaptation (LoRA), which freezes the pre-trained backbone weights $W_0 \in \mathbb{R}^{d \times k}$ and injects trainable rank-decomposition matrices $\Delta W = B \cdot A$, where $B \in \mathbb{R}^{d \times r}$, $A \in \mathbb{R}^{r \times k}$, and rank $r \ll \min(d, k)$. Beyond drastically reducing trainable parameters and GPU memory overhead, recent theoretical insights suggest that constraining updates to low-rank subspaces acts as implicit regularization, mitigating catastrophic forgetting of pre-trained linguistic priors. In this paper, we hypothesize that this regularization makes LoRA inherently more resilient to input orthographic shifts than full parameter updates.
